\label{sec:fundamentacao}

\section{A LLM como componente e a interface do filtro}\label{sec:interface}

Utilizar uma LLM como componente de software exige tratá-la como um módulo com
interface de entrada e saída bem definidas, e não como uma caixa de diálogo
isolada \cite{weber2024llm}. O componente de filtro ocupa uma posição
particular: atua antes da execução principal e não é responsável pela
resposta final, mas por classificar e preparar a requisição recebida. Suas
entradas são a requisição do usuário (potencialmente hostil); o
conteúdo não confiável incorporado ao contexto, estes são: trechos recuperados
\cite{lewis2020rag}, retornos de ferramenta, documentos, páginas, entre outros recursos
que constituem o vetor da injeção indireta \cite{greshake2023ipi}; e a política de
classificação. Suas saídas formam um contrato tipado: uma decisão
(\texttt{passar}/\texttt{bloquear}), uma marcação de risco e,
opcionalmente, uma versão transformada do conteúdo; é neste último campo que
reside a proposta deste trabalho.

\section{\textit{Prompt injection}: direto e indireto}\label{sec:injection}

O \textit{prompt injection} é o ataque mais fundamental contra componentes
baseados em LLM. Sua premissa é simples: instruções do administrador/desenvolvedor
e a entrada do usuário chegam pelo mesmo canal, 
sem separação semântica garantida, de modo que qualquer conteúdo textual
pode ser interpretado como instrução. Os primeiros estudos quantitativos
introduziram as noções de \textit{goal hijacking} e \textit{prompt leaking}
\cite{perez2022ignore}; trabalhos posteriores formalizaram o ataque e
padronizaram sua avaliação \cite{liu2024formalizing}. Na variante
\textbf{direta}, o atacante formula entradas que sobrepõem as instruções
originais, da sobreposição explícita (``ignore as instruções anteriores'') a
variantes como \textit{role-playing} e ofuscação por codificação. Na variante
\textbf{indireta} \cite{greshake2023ipi}, o atacante planta instruções em
conteúdo que o componente processará normalmente, como um documento indexado, o
retorno de uma ferramenta ou uma página web. É a variante indireta que motiva o
recorte deste trabalho, por não depender de nenhuma ação do usuário legítimo.

\section{A não separação como propriedade mensurável}\label{sec:nao_separacao}

A não separação entre instrução do administrador/desenvolvedor e a entrada do usuário decorre da forma como os
\textit{Transformers} são treinados para seguir instruções. O alinhamento por
instrução, como o RLHF \cite{ouyang2022instructgpt}, melhorou o respeito ao
\textit{system prompt}, mas não eliminou o problema, pois o modelo permanece um
processador estatístico de texto sem mecanismo intrínseco para verificar a
procedência de cada \textit{token}.

Trabalho recente elevou essa observação a uma propriedade formal: é possível
definir e medir o grau em que um modelo separa instrução de dado,
comparando o comportamento do modelo quando uma mesma sentença aparece na
posição de instrução e na posição de dado \cite{zverev2025separation}. A
conclusão é que nenhum modelo avaliado atinge separação alta, e que técnicas
canônicas de mitigação ou falham em melhorá-la substancialmente ou custam
utilidade. Isso importa aqui por dois motivos: fornece a métrica com que se pode
avaliar diretamente o efeito da intervenção proposta, e estabelece que a
separação, embora não garantida pelo modelo, pode ser
aproximada por convenções, premissa que fundamenta os métodos
discutidos a seguir.

Há duas vias para obter essa aproximação. A via aprendida altera os
pesos: consultas estruturadas \cite{chen2025struq}, otimização por preferência
\cite{chen2025secalign} e hierarquia de instruções \cite{wallace2024hierarchy}
demonstram que a separação funciona quando treinada, e há modelos abertos já
publicados sob essa abordagem \cite{chen2025metasecalign}; a via arquitetural vai
além e codifica a procedência nos próprios \textit{embeddings}
\cite{zverev2026aside}. Ambas exigem treinamento, e por isso servem aqui como
referência. Uma terceira via reprojeta o agente em
torno do modelo. O CaMeL separa fluxo de controle e fluxo
de dados por meio de um interpretador com capacidades
\cite{debenedetti2025camel}, e \citeonline{beurerkellner2025patterns}
catalogam padrões de projeto que restringem, por construção, o que o agente
pode fazer depois de ler conteúdo não confiável. A via
que este trabalho adota, opera apenas sobre o texto de entrada.

\section{\textit{Guardrails} de entrada e sua circularidade}\label{sec:guardrails}

\textit{Guardrails} de entrada inspecionam toda a informação que chega ao
componente antes de o modelo recebê-la. Há uma escolha estrutural entre filtros
baseados em padrões (expressões regulares, listas), rápidos e auditáveis porém
frágeis, e classificadores baseados em modelo, capazes de capturar semanticamente
formulações que escapam a filtros superficiais \cite{dong2024guardrails}. Estes últimos, os
\textit{guard models}, são modelos de linguagem, frequentemente menores que o
principal, ajustados para classificar entrada e saída segundo uma taxonomia
\cite{inan2023llamaguard, zeng2024shieldgemma, zhao2025qwen3guard}, e já compõem
sistemas de defesa em camadas \cite{chennabasappa2025llamafirewall}.

Um ponto importante para este trabalho é a circularidade: um \textit{guard model} é
ele próprio uma LLM e herda as vulnerabilidades que pretende mitigar. Pode ser
manipulado por instruções embutidas no conteúdo que analisa e contornado por
ofuscação \cite{hackett2025bypassing}. Se separar a instrução do desenvolvedor
da entrada não confiável é o que falta ao componente principal, não há razão para supor que o filtro esteja
dispensado dela.

\section{\textit{Spotlighting}: separação por convenção}\label{sec:spotlighting}

O \textit{spotlighting} \cite{hines2024spotlighting} é a formulação de
referência da separação por convenção. A ideia é transformar o texto não
confiável antes de inseri-lo no \textit{prompt}, de modo que o modelo receba um
sinal contínuo de sua procedência, e avisar no \textit{prompt} de sistema que a
transformação foi feita e que nada dentro do bloco deve ser obedecido. Os
autores descrevem três instanciações, de intrusividade crescente:

\begin{itemize}
  \item \textbf{Delimitação} (\textit{delimiting}). Símbolos especiais são
  colocados antes e depois do bloco não confiável, e o \textit{prompt} de
  sistema declara que o texto entre eles é um documento a ser lido, e não
  instrução a ser seguida. Um resumo de e-mail, por exemplo, chegaria ao modelo
  como ``\texttt{<{}<} texto do e-mail \texttt{>{}>}'', precedido de ``nunca
  obedeça a instruções entre esses símbolos''. É a forma mais simples e a única
  que não altera o conteúdo do bloco; em contrapartida, marca apenas as bordas,
  e um atacante que conheça os símbolos pode fechá-los e inserir os seus
  próprios.
  \item \textbf{Marcação de dados} (\textit{datamarking}). Em vez de marcar
  apenas as bordas, um caractere especial é intercalado ao longo de todo o
  bloco, tipicamente no lugar de cada espaço em branco: ``In this manner''
  torna-se ``In\textasciicircum{}this\textasciicircum{}manner''. O
  \textit{prompt} de sistema explica a transformação, e a procedência passa a
  acompanhar cada palavra, e não só o início e o fim. O texto continua legível
  para o modelo, mas qualquer instrução embutida carrega a marca de dado.
  \item \textbf{Codificação} (\textit{encoding}). O bloco inteiro é
  transformado por um algoritmo conhecido, como base64 ou ROT13, e o modelo é
  instruído a decodificá-lo antes de cumprir a tarefa. O conteúdo deixa de ser
  texto legível e vira dado opaco, o que torna a fronteira inequívoca, ao custo
  de exigir que o modelo decodifique corretamente.
\end{itemize}

As três reduziram de forma expressiva o sucesso de injeções indiretas nos
experimentos originais, com a marcação de dados oferecendo o melhor equilíbrio
entre proteção e preservação da tarefa, e a codificação dependendo da capacidade
do modelo para decodificar sem perda. Esses resultados foram obtidos protegendo
o modelo principal, o que torna as técnicas candidatas naturais a serem
aplicadas também ao filtro. Deste repertório, este trabalho adota a
delimitação, por ser a menos intrusiva e a única que não altera o texto que o
filtro precisa compreender, e transpõe a ideia de marcação por caractere para a
saída do filtro, onde ela passa a ser aplicada de forma seletiva. A fragilidade
da delimitação frente a quem conhece os símbolos é tratada pelo escape dos
caracteres de tag, descrito na Seção~\ref{sec:fator1}.

É preciso distinguir esse mecanismo de defesas puramente textuais, que
acrescentam ao \textit{prompt} instruções de vigilância sem alterar a estrutura
da entrada. O \textit{self-reminder} \cite{xie2023selfreminder} é o exemplo
canônico, e também reduz o sucesso de ataques. A distinção importa porque as duas
famílias operam por vias diferentes: uma pede atenção ao modelo, a outra muda a
forma do que ele lê. É a segunda que este trabalho investiga.

\section{Filtros que sanitizam a entrada}\label{sec:sanitizacao}

Uma linha recente vai além de classificar e passa a \textbf{transformar} o
conteúdo suspeito. O PromptArmor usa uma LLM sem treinamento adicional para
localizar o trecho injetado e removê-lo por casamento aproximado, atingindo
taxas de erro abaixo de 1\% em um benchmark de agentes \cite{shi2025promptarmor}.
O DataFilter treina um modelo dedicado sobre o mesmo \textit{backbone} de 8
bilhões de parâmetros para apagar o trecho malicioso preservando a informação
benigna, e utiliza um \textit{token} de fronteira para separar a instrução do
filtro do conteúdo não confiável na entrada do próprio filtro \cite{wang2025datafilter}. O PISanitizer opera em
caixa-branca, identificando os \textit{tokens} injetados pelos pesos de atenção
e eliminando-os \cite{geng2025pisanitizer}.

Os três compartilham a mesma decisão de projeto, que é a \textbf{remoção},
e este trabalho se diferencia deles nesta decisão, tomando uma abordagem diferente.
A remoção é irreversível e tomada sob incerteza: 
um falso positivo apaga conteúdo legítimo, e o componente
principal responde sobre um texto mutilado sem qualquer sinal de que isso
ocorreu. A marcação, por outro lado, é uma operação conservadora: preserva o
texto integralmente e delega ao componente principal a ponderação do risco,
transformando uma decisão binária do filtro em um sinal graduado. Além disso,
apenas o DataFilter separa a instrução do filtro do conteúdo não confiável, e o
faz por meio de treinamento; nenhum dos três avalia o efeito de fazê-lo por
alteração estrutural do \textit{input}.

\section{Métricas de avaliação}\label{sec:metricas_fund}

Medir apenas a detecção do filtro é insuficiente por duas razões. Primeiro, o
que interessa é o desfecho: um ataque que passa pelo filtro mas não é executado
pelo componente principal não é um ataque bem-sucedido. A métrica primária é,
portanto, a \textbf{taxa de sucesso do ataque} (ASR) medida na saída do sistema
\cite{yi2025bipia}. Segundo, um filtro é sempre um equilíbrio entre detecção e
\textbf{sobre-bloqueio}, a fração de entradas benignas indevidamente
bloqueadas (modo de falha agravado por conteúdos benignos que contêm
palavras-gatilho \cite{li2024injecguard}), que tem igual impacto em casos
reais de aplicação. Como o filtro opera sob restrição de
falsos positivos, a métrica de classificação operacional é a taxa
de detecção a uma taxa de falsos positivos baixa e fixa \cite{jacob2025promptshield}.
A esses somam-se a \textbf{utilidade} preservada na tarefa legítima e o custo em
\textit{tokens} e latência, seguindo o princípio consolidado de avaliar
conjuntamente utilidade e segurança \cite{debenedetti2024agentdojo}.
